\documentclass[11pt]{article}
\usepackage{fontspec}
\setmainfont[
  Path=fonts/,
  Extension=.ttf,
  UprightFont=*-400,
  BoldFont=*-700
]{NotoSerifSC}
\setsansfont[
  Path=fonts/,
  Extension=.ttf,
  UprightFont=*-400,
  BoldFont=*-700
]{NotoSerifSC}
\XeTeXlinebreaklocale "zh"
\XeTeXlinebreakskip=0pt plus 1pt
\usepackage[a4paper,margin=1.70cm]{geometry}
\usepackage{amsmath,amssymb,mathtools,bm}
\usepackage{booktabs,longtable,array}
\usepackage{enumitem}
\setlength{\parindent}{2em}
\setlength{\parskip}{0.35em}
\setlist{nosep,leftmargin=2em}
\allowdisplaybreaks[3]
\numberwithin{equation}{section}
\pagestyle{plain}

\newcommand{\R}{\mathbb{R}}
\newcommand{\E}{\mathbb{E}}
\newcommand{\Pp}{\mathbb{P}}
\newcommand{\1}{\mathbf{1}}
\newcommand{\T}{\mathsf{T}}
\newcommand{\F}{\mathrm{F}}
\newcommand{\cN}{\mathcal{N}}
\newcommand{\cL}{\mathcal{L}}
\newcommand{\cS}{\mathcal{S}}
\newcommand{\cH}{\mathcal{H}}
\newcommand{\cA}{\mathcal{A}}
\newcommand{\cM}{\mathcal{M}}
\newcommand{\cT}{\mathcal{T}}
\newcommand{\diag}{\operatorname{diag}}
\newcommand{\tr}{\operatorname{tr}}
\newcommand{\Var}{\operatorname{Var}}
\newcommand{\Cov}{\operatorname{Cov}}
\newcommand{\softmax}{\operatorname{softmax}}
\newcommand{\KL}{D_{\mathrm{KL}}}
\newcommand{\sg}{\operatorname{sg}}
\begin{document}
    \begin{center}
      {\LARGE\bfseries 5-1712.01815v1-Reinforcement-AlphaGo Zero}
    \end{center}
    \vspace{0.8em}

    \section{符号与维度}
    \begin{longtable}{>{$}l<{$} >{$}l<{$} p{10cm}}
    \toprule
    \text{符号} & \text{维度} & \text{含义}\\
    \midrule
    P,Q,N & -- & MCTS 先验、均值价值与访问计数。\\
        \pi,p_\theta & \Delta^{|\cA|-1} & 搜索策略与网络策略。\\
        v_\theta,z & [-1,1] & 价值预测与最终胜负。\\
        \tau & \R_{>0} & 访问计数温度。\\
    \bottomrule
    \end{longtable}

    所有向量默认是列向量；未特别说明时，期望同时对数据、噪声和采样时刻取值。下文只展开论文中承担方法逻辑的公式，并把所需假设写在推导发生的位置。

\section{PUCT 选择项的来源与尺度}

树中边 $(s,a)$ 维护访问次数 $N(s,a)$、价值和 $W(s,a)$、均值
\begin{equation}
 Q(s,a)=\frac{W(s,a)}{N(s,a)}
\end{equation}
以及网络先验 $P(s,a)$。PUCT 选择
\begin{equation}
 a^*=\arg\max_a\{Q(s,a)+U(s,a)\},
 \qquad
 U(s,a)=c_{\rm puct}P(s,a)
 \frac{\sqrt{\sum_bN(s,b)}}{1+N(s,a)}.
\end{equation}
当一条边未访问时分母为 $1$，先验大的动作更先被探索；固定父节点总访问数，
$U$ 随该边访问次数约按 $1/N(s,a)$ 衰减。若所有动作最终都被无限访问，探索项趋于零，
选择逐渐由 $Q$ 主导。

对访问次数作连续近似，若多个动作的有效上界在平衡时接近，
\begin{equation}
 Q_a+cP_a\frac{\sqrt{N}}{1+N_a}\approx C,
\end{equation}
可解出
\begin{equation}
 N_a\approx\frac{cP_a\sqrt N}{C-Q_a}-1.
\end{equation}
因此访问量同时依赖先验与价值差，不等价于只按策略概率采样。

\section{叶子价值的反向备份}

一次模拟到叶子得到当前玩家视角价值 $v\in[-1,1]$。若玩家每层交替，沿路径反向备份
\begin{equation}
 N(s_t,a_t)\leftarrow N(s_t,a_t)+1,
 \qquad
 W(s_t,a_t)\leftarrow W(s_t,a_t)+(-1)^{L-t}v.
\end{equation}
符号翻转来自零和对弈中双方价值相反。更新后的均值
\begin{equation}
 Q^+=\frac{W+z}{N+1}
 =Q+\frac{z-Q}{N+1}
\end{equation}
是在线样本均值递推。若叶子估计 $z$ 条件无偏、方差 $\sigma^2$，独立访问下
$Q$ 的方差按 $\sigma^2/N$ 衰减。

\section{访问次数形成改进策略}

搜索后策略目标为
\begin{equation}
 \pi(a\mid s)=\frac{N(s,a)^{1/\tau}}
 {\sum_bN(s,b)^{1/\tau}}.
\end{equation}
$\tau\to0$ 时集中到最大访问动作，$\tau=1$ 时按访问频率，
$\tau\to\infty$ 时趋于均匀。令 $z_a=\log N(s,a)$，则
\begin{equation}
 \pi(a\mid s)=\operatorname{softmax}(z_a/\tau),
\end{equation}
所以温度对分布熵的影响与普通 softmax 完全相同。

网络策略头 $p_\theta$ 拟合搜索策略的交叉熵
\begin{equation}
 \cL_p=-\sum_a\pi(a\mid s)\log p_\theta(a\mid s).
\end{equation}
若 logits 为 $l$，则
\begin{equation}
 \frac{\partial\cL_p}{\partial l_a}=p_\theta(a\mid s)-\pi(a\mid s).
\end{equation}
因此搜索不是不可微模块的反向传播对象；它产生一个经规划改进的监督分布。

\section{价值、策略与正则的联合损失}

单个自博弈样本 $(s,\pi,z)$ 的联合目标
\begin{equation}
 \cL(\theta)=(z-v_\theta(s))^2
 -\pi^T\log p_\theta(\cdot\mid s)
 +c\|\theta\|_2^2.
\end{equation}
对应梯度为
\begin{equation}
 \nabla\cL
 =2(v_\theta-z)\nabla v_\theta
 -\sum_a\pi_a\nabla\log p_{\theta,a}
 +2c\theta.
\end{equation}
价值标签 $z$ 是最终胜负，搜索叶子使用的是当时网络价值 $v$；
两者处在不同时间尺度。训练使新网络拟合实际结果和搜索策略，再由新网络驱动下一轮搜索，
形成近似广义策略迭代。

\section{根节点狄利克雷噪声}

探索时根先验常替换为
\begin{equation}
 P'(s,a)=(1-\varepsilon)P(s,a)+\varepsilon\eta_a,
 \qquad \eta\sim\operatorname{Dir}(\alpha,\ldots,\alpha).
\end{equation}
对 $K$ 个动作，对称 Dirichlet 的矩为
\begin{equation}
 \E\eta_a=\frac1K,\qquad
 \Var(\eta_a)=\frac{K-1}{K^2(K\alpha+1)}.
\end{equation}
$\alpha$ 小时样本更尖锐，鼓励少数随机动作；$\alpha$ 大时接近均匀。
混合后
\begin{equation}
 \E P'_a=(1-\varepsilon)P_a+\frac\varepsilon K,
 \qquad
 \Var(P'_a)=\varepsilon^2\Var(\eta_a).
\end{equation}
噪声仅在根加入，避免整棵树每层都被独立噪声淹没。

\section{Bellman 期望方程与优势恒等式}

对策略 $\pi$，定义
\begin{align}
 V^\pi(s)&=\E_\pi\left[\sum_{t=0}^{\infty}\gamma^tr_t\mid s_0=s\right],\\
 Q^\pi(s,a)&=\E_\pi\left[\sum_{t=0}^{\infty}\gamma^tr_t
 \mid s_0=s,a_0=a\right].
\end{align}
分离首步奖励：
\begin{align}
 Q^\pi(s,a)&=r(s,a)+\gamma\E_{s'}V^\pi(s'),\\
 V^\pi(s)&=\E_{a\sim\pi(\cdot\mid s)}Q^\pi(s,a).
\end{align}
优势
\begin{equation}
 A^\pi(s,a)=Q^\pi(s,a)-V^\pi(s)
\end{equation}
在策略动作分布下均值为零：
\begin{equation}
 \E_{a\sim\pi}A^\pi(s,a)=0.
\end{equation}
这不表示每个动作优势小，只表示正负改进机会按当前策略加权抵消。

\section{性能差异引理}

令新旧策略为 $\pi',\pi$。沿 $\pi'$ 轨迹，
\begin{align}
 \E_{\pi'}\sum_{t=0}^{\infty}\gamma^tA^\pi(s_t,a_t)
 &=\E_{\pi'}\sum_t\gamma^t
 [r_t+\gamma V^\pi(s_{t+1})-V^\pi(s_t)]\\
 &=J(\pi')-\E V^\pi(s_0),
\end{align}
其中价值项望远镜消去，且有界价值使末项
$\gamma^{T+1}V(s_{T+1})\to0$。若初始分布相同，
$\E V^\pi(s_0)=J(\pi)$，故
\begin{equation}
 \boxed{J(\pi')-J(\pi)
 =\frac1{1-\gamma}
 \E_{s\sim d_{\pi'},a\sim\pi'}A^\pi(s,a).}
\end{equation}
策略优化代理常把难采样的 $d_{\pi'}$ 替换为旧占用分布 $d_\pi$；
信赖域或比率裁剪的目标是限制这一步替换造成的误差。

\section{KL 信赖域的二阶局部形式}

对小参数变化 $\Delta\theta$，
\begin{equation}
 D_{\rm KL}(\pi_{\theta}\|\pi_{\theta+\Delta})
 =\frac12\Delta^TF(\theta)\Delta+O(\|\Delta\|^3),
\end{equation}
其中 Fisher 信息
\begin{equation}
 F(\theta)=\E_{s,a}
 [\nabla\log\pi_\theta(a\mid s)
 \nabla\log\pi_\theta(a\mid s)^T].
\end{equation}
一阶目标改进 $g^T\Delta$，约束
$\frac12\Delta^TF\Delta\le\delta$。拉格朗日法：
\begin{equation}
 \mathcal J=g^T\Delta-\lambda
 \left(\frac12\Delta^TF\Delta-\delta\right).
\end{equation}
驻点
\begin{equation}
 g-\lambda F\Delta=0
 \Rightarrow \Delta=\lambda^{-1}F^{-1}g.
\end{equation}
代回约束得到
\begin{equation}
 \boxed{\Delta=\sqrt{\frac{2\delta}{g^TF^{-1}g}}F^{-1}g.}
\end{equation}
自然梯度 $F^{-1}g$ 是概率分布几何下最陡方向；PPO 裁剪是更易实现但不严格等价的近似。

\section{离策略序列比率的方差}

轨迹前 $T$ 步重要性比率
\begin{equation}
 W_T=\prod_{t=0}^{T-1}
 \frac{\pi(a_t\mid s_t)}{\mu(a_t\mid s_t)}.
\end{equation}
在行为策略 $\mu$ 下若支撑相容，$\E_\mu W_T=1$。
对数比率是和
\begin{equation}
 \log W_T=\sum_t\ell_t,\qquad
 \ell_t=\log\pi(a_t\mid s_t)-\log\mu(a_t\mid s_t).
\end{equation}
若粗略假设 $\ell_t$ 独立、均值 $m$、方差 $v$，则
$W_T$ 近似 log-normal，
\begin{equation}
 \frac{\Var(W_T)}{(\E W_T)^2}\approx e^{Tv}-1.
\end{equation}
方差随长度指数增长，这解释了 token/step 级截断、V-trace 或短视野校正的必要性。

截断 $\bar w_t=\min(c,w_t)$ 产生偏差，但单步权重有界，
能控制乘积和梯度爆炸。偏差精确来自尾部
\begin{equation}
 \E_\mu[(w_t-c)_+f_t].
\end{equation}

\section{熵正则的策略梯度}

状态熵
\begin{equation}
 \mathcal H(\pi_\theta)=-\sum_a\pi_\theta(a)\log\pi_\theta(a).
\end{equation}
求导
\begin{align}
 \nabla\mathcal H
 &=-\sum_a\nabla\pi(a)[1+\log\pi(a)]\\
 &=-\E_{a\sim\pi}
 [(1+\log\pi(a))\nabla\log\pi(a)].
\end{align}
因 $\E\nabla\log\pi=0$，常数 $1$ 可去掉：
\begin{equation}
 \nabla\mathcal H
 =-\E[\log\pi(a)\nabla\log\pi(a)].
\end{equation}
加入 $+\alpha\mathcal H$ 会提高低概率动作的相对激励，防止策略过早坍缩；
但在可验证答案任务中，过强熵也会降低正确答案集中度。

\section{奖励平移、缩放与归一化}

若所有奖励加常数 $c$，无限折扣回报
\begin{equation}
 G_t'=G_t+\frac{c}{1-\gamma}.
\end{equation}
精确优势不变，因为 $Q,V$ 同时加 $c/(1-\gamma)$。用采样回报而无正确基线时，
常数会增加梯度方差但期望 score 项仍抵消。

奖励乘正数 $a$ 时，优势和策略梯度乘 $a$，等价于改变有效学习率；
PPO 的裁剪边界不随 $a$ 变，所以裁剪区域不变但目标梯度尺度变。
组内标准化
\begin{equation}
 \widehat A_i=\frac{R_i-\bar R}{s_R+\epsilon}
\end{equation}
同时消除平移并近似消除尺度，但 $s_R$ 很小时会放大噪声，故需要 $\epsilon$ 或跳过全同奖励组。

\section{KL 正则与奖励塑形}

目标
\begin{equation}
 J(\pi)=\E_\pi[R]-\beta D_{\rm KL}(\pi\|\pi_{\rm ref})
\end{equation}
可写为每步塑形奖励
\begin{equation}
 r_t'=r_t-\beta
 \left[\log\pi(a_t\mid s_t)-\log\pi_{\rm ref}(a_t\mid s_t)\right].
\end{equation}
固定状态、给定动作价值 $Q(a)$ 时，最优非参数策略由拉格朗日法得到
\begin{equation}
 \pi^*(a)=\frac{\pi_{\rm ref}(a)e^{Q(a)/\beta}}
 {\sum_b\pi_{\rm ref}(b)e^{Q(b)/\beta}}.
\end{equation}
$\beta\to0$ 时集中到高价值动作，$\beta\to\infty$ 时回到参考策略。
参考策略为零概率的动作在正向 KL 约束下永远不可选，故支撑选择很重要。

\section{拒答动作与错误成本}

设回答正确收益 $u_c$、错误收益 $u_e$、拒答收益 $u_a$，模型置信度
$p=\Pp(\text{correct}\mid s)$。回答期望效用
\begin{equation}
 U_{\rm ans}(p)=pu_c+(1-p)u_e,
\end{equation}
拒答效用 $U_{\rm abstain}=u_a$。当
\begin{equation}
 p\ge\frac{u_a-u_e}{u_c-u_e}
\end{equation}
时回答更优（假设 $u_c>u_e$）。因此最优置信阈值由错误和拒答的相对成本决定，
不是固定的 $1/2$。RL 后训练若隐含奖励给拒答过低分，就会系统性鼓励猜测。

\section{有限 horizon 的回报界}

若 $|r_t|\le R_{\max}$，无限折扣回报有界
\begin{equation}
 |G_t|\le\sum_{k=0}^{\infty}\gamma^kR_{\max}
 =\frac{R_{\max}}{1-\gamma}.
\end{equation}
截断到 $H$ 步的尾部误差
\begin{equation}
 \left|G_t-\sum_{k=0}^{H-1}\gamma^kr_{t+k}\right|
 \le\frac{\gamma^HR_{\max}}{1-\gamma}.
\end{equation}
要使其不超过 $\epsilon$，充分条件
\begin{equation}
 H\ge\frac{\log[\epsilon(1-\gamma)/R_{\max}]}{\log\gamma},
\end{equation}
注意 $\log\gamma<0$，除法会反转不等号，右端为正。

\section{光滑损失的下降引理}

若 $L$ 的梯度是 $\beta$-Lipschitz：
\begin{equation}
 \|\nabla L(x)-\nabla L(y)\|
 \le\beta\|x-y\|,
\end{equation}
则沿线段积分
\begin{align}
 L(y)-L(x)
 &=\int_0^1\nabla L(x+t(y-x))^T(y-x)dt\\
 &=\nabla L(x)^T(y-x)+R,
\end{align}
余项满足
\begin{align}
 |R|
 &\le\int_0^1\beta t\|y-x\|^2dt
 =\frac\beta2\|y-x\|^2.
\end{align}
故下降引理
\begin{equation}
 L(y)\le L(x)+\nabla L(x)^T(y-x)
 +\frac\beta2\|y-x\|^2.
\end{equation}
取梯度步 $y=x-\eta\nabla L(x)$：
\begin{equation}
 L(y)\le L(x)-\eta\left(1-\frac{\beta\eta}{2}\right)
 \|\nabla L(x)\|^2.
\end{equation}
当 $0<\eta<2/\beta$ 时每个精确梯度步下降；深网中 $\beta$ 未知且随位置变化，
这仍解释了学习率过大造成不稳定的局部机制。

\section{二阶近似与曲率方向}

在 $x$ 附近 Taylor 展开
\begin{equation}
 L(x+\Delta)=L(x)+g^T\Delta
 +\frac12\Delta^TH\Delta+O(\|\Delta\|^3).
\end{equation}
若 $H=U\diag(\lambda_i)U^T$，在本征方向 $u_i$ 上步长 $a$ 的变化
\begin{equation}
 \Delta L\approx a g_i+\frac12\lambda_i a^2.
\end{equation}
$\lambda_i>0$ 为局部碗形，$\lambda_i<0$ 是负曲率逃离鞍点方向。
加阻尼的 Newton 步
\begin{equation}
 \Delta=-(H+\lambda I)^{-1}g
\end{equation}
在特征方向缩放为 $-g_i/(\lambda_i+\lambda)$，避免接近零或负特征值造成巨大步长。

\section{链式 Jacobian 与条件数}

复合映射 $y=f_L\circ\cdots\circ f_1(x)$ 的 Jacobian
\begin{equation}
 J_{y\leftarrow x}=J_LJ_{L-1}\cdots J_1.
\end{equation}
谱范数上界
\begin{equation}
 \|J_{y\leftarrow x}\|_2\le\prod_{\ell=1}^{L}\|J_\ell\|_2.
\end{equation}
若各层最小奇异值存在，
\begin{equation}
 \sigma_{\min}(J_{y\leftarrow x})
 \ge\prod_\ell\sigma_{\min}(J_\ell).
\end{equation}
因此条件数至多乘法累积：
\begin{equation}
 \kappa(J_{y\leftarrow x})
 \le\prod_\ell\kappa(J_\ell).
\end{equation}
低维映射、递归模块和物理传播都面临同一问题：表示存在不代表反向梯度数值良态。

\section{随机梯度的偏差与方差}

设样本梯度 $g(\theta;\xi)$ 满足
\begin{equation}
 \E_\xi g(\theta;\xi)=\nabla L(\theta),
 \qquad \E\|g-\nabla L\|^2\le\sigma^2.
\end{equation}
独立批量 $B$ 的平均
\begin{equation}
 \bar g_B=\frac1B\sum_{i=1}^{B}g_i
\end{equation}
满足
\begin{equation}
 \E\bar g_B=\nabla L,\qquad
 \E\|\bar g_B-\nabla L\|^2\le\frac{\sigma^2}{B}.
\end{equation}
若样本相关，协方差项出现：
\begin{equation}
 \Var(\bar g_B)=\frac1{B^2}\left[
 \sum_i\Var(g_i)+\sum_{i\ne j}\Cov(g_i,g_j)
 \right].
\end{equation}
扩大批量只有在交叉协方差不主导时才近似按 $1/B$ 降噪。

若使用有偏估计 $\E\widehat g=\nabla L+b$，均方误差分解
\begin{equation}
 \E\|\widehat g-\nabla L\|^2
 =\|b\|^2+\tr\Cov(\widehat g).
\end{equation}
直通估计、停止梯度目标和截断反传通常以可控偏差换低方差或低内存。

\section{Adam 的尺度归一化}

逐坐标 Adam 更新
\begin{align}
 m_t&=\beta_1m_{t-1}+(1-\beta_1)g_t,\\
 v_t&=\beta_2v_{t-1}+(1-\beta_2)g_t^2,\\
 \widehat m_t&=m_t/(1-\beta_1^t),\\
 \widehat v_t&=v_t/(1-\beta_2^t),\\
 \theta_{t+1}&=\theta_t-\eta
 \frac{\widehat m_t}{\sqrt{\widehat v_t}+\epsilon}.
\end{align}
若梯度长期乘常数 $c>0$，忽略 $\epsilon$ 时
$\widehat m$ 乘 $c$、$\sqrt{\widehat v}$ 也乘 $c$，更新近似不变。
但瞬态、符号变化和 $\epsilon$ 会破坏精确尺度不变性。

AdamW 把权重衰减解耦：
\begin{equation}
 \theta_{t+1}=(1-\eta\lambda)\theta_t
 -\eta\frac{\widehat m_t}{\sqrt{\widehat v_t}+\epsilon}.
\end{equation}
这不同于把 $\lambda\theta$ 加入梯度后再经过自适应分母；后者会按坐标二阶矩缩放正则化。

\section{范数约束与投影梯度}

约束集合 $\mathcal C$ 上的投影梯度
\begin{equation}
 \theta^+=\Pi_{\mathcal C}(\theta-\eta g),
 \qquad
 \Pi_{\mathcal C}(x)=\arg\min_{y\in\mathcal C}\|x-y\|^2.
\end{equation}
对欧氏球 $\mathcal C=\{x:\|x\|\le R\}$，
\begin{equation}
 \Pi_{\mathcal C}(x)=
 \begin{cases}
 x,&\|x\|\le R,\\
 Rx/\|x\|,&\|x\|>R.
 \end{cases}
\end{equation}
对流形约束，局部一阶步先投影到切空间，再用 retraction 回到流形；
仅把参数写成 $g(z)$ 相当于用参数化隐式满足约束，但会引入 $J_g^TJ_g$ 的坐标度量。

\section{Lipschitz 误差的层间传播}

真实模块 $f_\ell$、近似模块 $\widehat f_\ell$ 满足
\begin{equation}
 \|f_\ell(x)-\widehat f_\ell(x)\|\le\epsilon_\ell,
 \qquad
 \|f_\ell(x)-f_\ell(y)\|\le L_\ell\|x-y\|.
\end{equation}
令真实、近似中间状态差 $e_\ell$。三角不等式给出
\begin{align}
 e_{\ell+1}
 &=\|f_\ell(h_\ell)-\widehat f_\ell(\widehat h_\ell)\|\\
 &\le L_\ell e_\ell+\epsilon_\ell.
\end{align}
递归展开
\begin{equation}
 e_L\le
 \left(\prod_{j=0}^{L-1}L_j\right)e_0
 +\sum_{\ell=0}^{L-1}
 \left(\prod_{j=\ell+1}^{L-1}L_j\right)\epsilon_\ell.
\end{equation}
早期层误差乘更多后续 Lipschitz 因子，所以量化、低秩近似或局部训练的相同单层误差，
放在不同深度并不等价。

\section{复杂度必须区分四个量}

参数量、FLOPs、激活内存和通信量是不同指标。以线性层
$Y=XW$、$X\in\R^{B\times n}$、$W\in\R^{n\times m}$ 为例：
\begin{align}
 \text{参数量}&=nm,\\
 \text{前向乘加量}&\asymp Bnm,\\
 \text{主要激活量}&=B(n+m),\\
 \text{数据并行梯度通信}&\asymp nm.
\end{align}
低秩因子可同时减少参数与乘法，但动态稀疏若实现仍做稠密 kernel，参数/FLOPs 公式不会自动转化为墙钟加速。
冻结参数减少梯度和优化器状态，却不一定减少前向计算或输入梯度计算。

\end{document}
